\documentclass[a4paper,12pt]{article}
\usepackage[T2A]{fontenc}
\usepackage[utf8]{inputenc}
\usepackage[english,russian]{babel}
\usepackage{amsmath}
\usepackage{amssymb}
\usepackage{bm} 

\begin{document}

\title{Лабораторная работа №3 \\ Математический режим LaTeX}
\author{Сунь Маосин}
\date{2026}
\maketitle

\section{Цель работы}

Изучение математического режима LaTeX: inline и display формулы, индексы, греческие буквы, матрицы и выравнивание с помощью пакета amsmath.

\section{Math mode (математический режим)}

\subsection{Inline math (внутритекстовые формулы)}

Формула внутри текста: $E = mc^2$. 
Ещё пример: $a^2 + b^2 = c^2$.

Можно использовать скобки: \( \sin^2 x + \cos^2 x = 1 \).

\subsection{Display math (выключные формулы)}

Формула на отдельной строке:
\[
E = mc^2
\]

Ещё пример:
\[
\int_{0}^{\infty} e^{-x^2} dx = \frac{\sqrt{\pi}}{2}
\]

Нумерованная формула:
\begin{equation}
\sum_{n=1}^{\infty} \frac{1}{n^2} = \frac{\pi^2}{6}
\end{equation}

\section{Индексы и греческие буквы}

\subsection{Верхние и нижние индексы}

$x^2$ — верхний индекс \\
$x_2$ — нижний индекс \\
$x^{2y}$ — сложный верхний индекс \\
$x_{ij}$ — сложный нижний индекс \\
$x^{2}_{i}$ — комбинация

\subsection{Греческие буквы}

Строчные: $\alpha, \beta, \gamma, \delta, \epsilon, \varepsilon, \zeta, \eta, \theta, \iota, \kappa, \lambda, \mu, \nu, \xi, \pi, \rho, \sigma, \tau, \upsilon, \phi, \varphi, \chi, \psi, \omega$

Заглавные: $\Gamma, \Delta, \Theta, \Lambda, \Xi, \Pi, \Sigma, \Upsilon, \Phi, \Psi, \Omega$

\section{Математические функции}

$\sin x$, $\cos x$, $\tan x$, $\cot x$ — тригонометрические функции \\
$\arcsin x$, $\arccos x$, $\arctan x$ — обратные функции \\
$\log x$, $\ln x$, $\lg x$ — логарифмы \\
$\lim_{x \to 0} \frac{\sin x}{x} = 1$ — предел \\
$\max_{x \in [0,1]} f(x)$ — максимум \\
$\min_{x \in [0,1]} f(x)$ — минимум

\section{Дроби, корни и биномы}

Обычная дробь: $\frac{1}{2}$ \\
Сложная дробь: $\frac{x^2 + y^2}{x - y}$ \\
Квадратный корень: $\sqrt{2}$ \\
Корень n-ной степени: $\sqrt[n]{x}$ \\
Биномиальные коэффициенты: $\binom{n}{k} = \frac{n!}{k!(n-k)!}$

\section{Интегралы, суммы и произведения}

Определённый интеграл: $\int_{0}^{1} x^2 \, dx$ \\
Двойной интеграл: $\iint_{D} f(x,y) \, dx \, dy$ \\
Тройной интеграл: $\iiint_{V} f(x,y,z) \, dx \, dy \, dz$ \\
Сумма: $\sum_{n=1}^{\infty} \frac{1}{n^2}$ \\
Произведение: $\prod_{i=1}^{n} i = n!$

\section{Матрицы}

Матрица без скобок:
\[
\begin{matrix}
a & b \\
c & d
\end{matrix}
\]

Матрица в круглых скобках:
\[
\begin{pmatrix}
a & b \\
c & d
\end{pmatrix}
\]

Матрица в квадратных скобках:
\[
\begin{bmatrix}
1 & 2 & 3 \\
4 & 5 & 6 \\
7 & 8 & 9
\end{bmatrix}
\]

Матрица в вертикальных линиях (определитель):
\[
\begin{vmatrix}
a & b \\
c & d
\end{vmatrix} = ad - bc
\]

\section{Многострочные формулы с amsmath}

\subsection{Системы уравнений}

\[
\begin{cases}
x + y = 10 \\
x - y = 4
\end{cases}
\]

\subsection{Выравнивание с align}

\begin{align}
x + y &= 10 \\
x - y &= 4 \\
2x &= 14 \\
x &= 7
\end{align}

\subsection{Выравнивание без номеров}

\begin{align*}
f(x) &= x^2 + 2x + 1 \\
     &= (x + 1)^2
\end{align*}

\subsection{Несколько столбцов}

\begin{align*}
x &= y & z &= w & a &= b \\
x^2 &= y^2 & z^2 &= w^2 & a^2 &= b^2
\end{align*}

\subsection{Длинные формулы с переносом}

\begin{multline}
x = a + b + c + d + e + f + g + h + i + j \\
+ k + l + m + n + o + p + q + r + s + t + u
\end{multline}

\subsection{Группировка формул}

\begin{gather}
x + y = 10 \\
x - y = 4 \\
2x = 14
\end{gather}

\section{Математические шрифты}

$\mathrm{Roman}$ — прямой шрифт \\
$\mathit{Italic}$ — курсив (для слов) \\
$\mathbf{Bold}$ — жирный \\
$\mathsf{Sans\ serif}$ — рубленый \\
$\mathtt{Typewriter}$ — пишущая машинка \\
$\mathcal{CALIGRAPHIC}$ — каллиграфический \\
$\mathbb{BLACKBOARD\ BOLD}$ — двойной (для множеств) \\
$\mathfrak{Fraktur}$ — готический

\section{Жирные символы с bm}

$\bm{\alpha} + \bm{\beta}$ — жирные греческие \\
$\bm{=}$ — жирный знак равенства \\
$\bm{\nabla}$ — жирный набла

\section{Текст внутри математики}

\[
\text{Область определения: } \{x \in \mathbb{R} \mid x > 0\}
\]

\[
\mathbb{N} \subset \mathbb{Z} \subset \mathbb{Q} \subset \mathbb{R} \subset \mathbb{C}
\]

\section{Вывод}

В ходе работы были изучены основные возможности математического режима LaTeX:
\begin{itemize}
  \item inline и display формулы
  \item индексы и греческие буквы
  \item дроби, корни, интегралы и суммы
  \item матрицы различных типов
  \item многострочные формулы с пакетом amsmath
  \item математические шрифты
\end{itemize}

\end{document}